\subsection{Modules over Algebras}

\begin{definition} \label{def:subalgebra-generated-subalgebra} \label{def:algebra-ideal-simple-quotient} \label{def:algebra-homomorphism} \label{def:left-regular-representation} \hlt{Subalgebras, Ideals, Quotients, Homomorphisms, and Regular Representation}\\
Let $\mathcal{A}$ be an $\mathbb{F}$-algebra.
\begin{enumerate}[label=(\roman*)]
\setlength{\itemsep}{0pt}
\item A \hlt{subalgebra} is a subset $\mathcal{B}\subseteq\mathcal{A}$ such that $1_{\mathcal{A}}\in\mathcal{B}$, $\mathcal{B}\leq_{\mathbb{F}}\mathcal{A}$, and $b,b'\in\mathcal{B}\Rightarrow bb'\in\mathcal{B}$.
\item For $X\subseteq\mathcal{A}$, the \hlt{subalgebra generated by $X$} is $\langle X\rangle_{\operatorname{subalg}}\defeq\bigcap\{\mathcal{B}:\mathcal{B}\text{ is a subalgebra of }\mathcal{A},\ X\subseteq\mathcal{B}\}$.\\
Equivalently, $\langle X\rangle_{\operatorname{subalg}}=\operatorname{span}_{\mathbb{F}}\{1_{\mathcal{A}},\ x_1\cdots x_m:m\geq1,\ x_i\in X\}$.
\item A \hlt{two-sided ideal} is a subset $\mathcal{I}\subseteq\mathcal{A}$ such that $\mathcal{I}\neq\emptyset$, $u,v\in\mathcal{I}\Rightarrow u-v\in\mathcal{I}$, and $a\in\mathcal{A},\ u\in\mathcal{I}\Rightarrow au,ua\in\mathcal{I}$.
\item For nonempty $X\subseteq\mathcal{A}$, the \hlt{ideal generated by $X$} is $\langle X\rangle_{\operatorname{ideal}}\defeq\{\sum_{j=1}^{n}a_jx_jb_j:n\geq1,\ x_j\in X,\ a_j,b_j\in\mathcal{A}\}$.
\item $\mathcal{A}$ is \hlt{simple} iff $\mathcal{A}^2\neq0$ and its only two-sided ideals are $0$ and $\mathcal{A}$.
\item If $\mathcal{I}\triangleleft\mathcal{A}$, then $\mathcal{A}/\mathcal{I}$ is an $\mathbb{F}$-algebra by $(a+\mathcal{I})(b+\mathcal{I})\defeq ab+\mathcal{I}$ and $r(a+\mathcal{I})\defeq ra+\mathcal{I}$.
\item For $\mathbb{F}$-algebras $\mathcal{A},\mathcal{B}$, an \hlt{algebra homomorphism} is a map $\sigma:\mathcal{A}\to\mathcal{B}$ such that, for all $a,a'\in\mathcal{A}$ and $r\in\mathbb{F}$,\\
$\sigma(a+a')=\sigma(a)+\sigma(a')$, $\sigma(aa')=\sigma(a)\sigma(a')$, $\sigma(1_{\mathcal{A}})=1_{\mathcal{B}}$, and $\sigma(ra)=r\sigma(a)$.
\item For $a\in\mathcal{A}$, the \hlt{left multiplication operator} is $L_a:\mathcal{A}\to\mathcal{A}$, $x\mapsto ax$.\\
The map $L_{\mathcal{A}}:\mathcal{A}\to\LL(\mathcal{A})$, $a\mapsto L_a$, is the \hlt{left regular representation} of $\mathcal{A}$.
\end{enumerate}
\end{definition}

\begin{proposition} \label{prop:algebra-kernel-first-isomorphism} \hlt{Kernel and First Isomorphism Theorem for Algebras}\\
If $\sigma:\mathcal{A}\to\mathcal{B}$ is an $\mathbb{F}$-algebra homomorphism,\\
then $\ker\sigma\triangleleft\mathcal{A}$ and $\overline{\sigma}:\mathcal{A}/\ker\sigma\to\operatorname{im}\sigma$, $a+\ker\sigma\mapsto\sigma(a)$, is an $\mathbb{F}$-algebra isomorphism.
\begin{proof} Set $K\defeq\ker\sigma$. Linearity gives $K\leq_{\mathbb{F}}\mathcal{A}$.\\
If $u\in K$ and $a\in\mathcal{A}$, then $\sigma(au)=\sigma(a)\sigma(u)=0=\sigma(u)\sigma(a)=\sigma(ua)$, hence $au,ua\in K$ and $K\triangleleft\mathcal{A}$.\\
Define $\overline{\sigma}(a+K)\defeq\sigma(a)$. Since $a+K=a'+K\Longleftrightarrow a-a'\in K\Longleftrightarrow\sigma(a)=\sigma(a')$, $\overline{\sigma}$ is well defined and injective; it is surjective onto $\operatorname{im}\sigma$.\\
For $a,a'\in\mathcal{A}$ and $r\in\mathbb{F}$, $\overline{\sigma}((a+K)(a'+K))=\sigma(aa')=\sigma(a)\sigma(a')$, $\overline{\sigma}(r(a+K))=r\sigma(a)$, and $\overline{\sigma}(1_{\mathcal{A}}+K)=1_{\mathcal{B}}$.\\
Thus $\overline{\sigma}$ is an $\mathbb{F}$-algebra isomorphism.\\
\hlt{Structural Commentary}: Quotienting by the kernel removes exactly the relations invisible to $\sigma$.\\
\hlt{Key Insight}: The image is the faithful algebra obtained after collapsing $\ker\sigma$.
\end{proof}
\end{proposition}

\begin{theorem} \label{thm:left-regular-representation-algebra} \hlt{Left Regular Representation}\\
For every $\mathbb{F}$-algebra $\mathcal{A}$, $L_{\mathcal{A}}:\mathcal{A}\hookrightarrow\LL(\mathcal{A})$, $a\mapsto L_a$, is an $\mathbb{F}$-algebra monomorphism.
\begin{proof}
For $a,b\in\mathcal{A}$, $r\in\mathbb{F}$, and $x\in\mathcal{A}$, $L_{a+b}(x)=(a+b)x=L_a(x)+L_b(x)$ and $L_{ra}(x)=(ra)x=rL_a(x)$, so $L_{\mathcal{A}}$ is $\mathbb{F}$-linear. Also $L_{ab}(x)=(ab)x=a(bx)=(L_aL_b)(x)$ and $L_{1_{\mathcal{A}}}=I_{\mathcal{A}}$, so products and units are preserved.\\
If $L_a=0$, then $a=L_a(1_{\mathcal{A}})=0$, so $L_{\mathcal{A}}$ is injective.\\
\hlt{Structural Commentary}: Left multiplication realises $\mathcal{A}$ inside its own operator algebra.\\
\hlt{Key Insight}: The unit detects elements, since $L_a(1_{\mathcal{A}})=a$.
\end{proof}
\end{theorem}

\begin{figure}[H]
  \centering
  \begin{tikzpicture}[
  >=Stealth,
  x=0.9cm,
  every node/.style={font=\small},
  alg/.style={draw=blue!60!black, fill=blue!7, rounded corners=2pt,
    minimum width=2.2cm, minimum height=0.64cm, align=center, inner sep=4pt},
  op/.style={draw=dartmouthgreen!65!black, fill=dartmouthgreen!8, rounded corners=2pt,
    minimum width=2.65cm, minimum height=0.64cm, align=center, inner sep=4pt},
  vec/.style={draw=gray!55, fill=gray!6, rounded corners=2pt,
    minimum width=1.8cm, minimum height=0.56cm, align=center, inner sep=3pt},
  arr/.style={->, line width=0.75pt, draw=black!70},
  note/.style={font=\scriptsize, inner sep=2pt, fill=white}
]
  \node[alg] (A) at (0,0) {$a\in\mathcal{A}$};
  \node[op] (L) at (4.15,0) {$L_a\in\LL(\mathcal{A})$};
  \node[vec] (x) at (8.05,0) {$x\in\mathcal{A}$};
  \node[vec] (ax) at (11.45,0) {$ax\in\mathcal{A}$};

  \draw[arr] (A) -- node[note, above] {$a\mapsto L_a$} (L);
  \draw[arr] (L) -- node[note, above] {acts on} (x);
  \draw[arr] (x) -- node[note, above] {$L_a$} (ax);
\end{tikzpicture}

  \caption{The left regular representation sends $a\in\mathcal{A}$ to
  $L_a:x\mapsto ax$; $L_{ab}=L_aL_b$ and $L_a(1_{\mathcal{A}})=a$ give
  multiplicativity and injectivity.}
  \label{fig:algebra-left-regular-representation}
\end{figure}

\begin{theorem} \label{thm:endomorphism-algebra-center-ideals-simple} \hlt{Center and Ideals of an Endomorphism Algebra}\\
Let $0\neq\mathcal{V}$ be an $\mathbb{F}$-vector space, and set $\mathcal{F}(\mathcal{V})\defeq\{T\in\LL(\mathcal{V}):\dim\operatorname{im}T<\infty\}$.
\begin{enumerate}[label=(\roman*)]
\setlength{\itemsep}{0pt}
\item $Z(\LL(\mathcal{V}))=\mathbb{F}I_{\mathcal{V}}$, so $\LL(\mathcal{V})$ is central.
\item $\mathcal{F}(\mathcal{V})\triangleleft\LL(\mathcal{V})$, and $0\neq\mathcal{J}\triangleleft\LL(\mathcal{V})\Rightarrow\mathcal{F}(\mathcal{V})\subseteq\mathcal{J}$.
\item $\LL(\mathcal{V})$ is simple iff $\dim\mathcal{V}<\infty$.
\end{enumerate}
\begin{proof}
\begin{enumerate}[label=(\roman*)]
\setlength{\itemsep}{0pt}
\item Let $T\in Z(\LL(\mathcal{V}))$. For $0\neq v\in\mathcal{V}$ and $\varphi\in\mathcal{V}^*$ with $\varphi(v)=1$, commutation with $v\otimes\varphi:x\mapsto\varphi(x)v$ gives $Tv=T(v\otimes\varphi)(v)=(v\otimes\varphi)T(v)=\varphi(Tv)v$, so every line is $T$-stable. Thus $Tv=\lambda_vv$ for $v\neq0$, and applying this to independent $u,v$ and to $u+v$ gives $\lambda_u=\lambda_v$; hence $T=rI_{\mathcal{V}}$.\\
Scalars commute with all endomorphisms.
\item If $R,S\in\LL(\mathcal{V})$ and $T\in\mathcal{F}(\mathcal{V})$, then $\operatorname{im}(RTS)\subseteq R(\operatorname{im}T)$, so $RTS\in\mathcal{F}(\mathcal{V})$ and $\mathcal{F}(\mathcal{V})\triangleleft\LL(\mathcal{V})$. If $0\neq\mathcal{J}\triangleleft\LL(\mathcal{V})$, choose $0\neq S\in\mathcal{J}$, $e$ with $Se\neq0$, and $\eta\in\mathcal{V}^*$ with $\eta(Se)=1$. For any $v\in\mathcal{V}$ and $\varphi\in\mathcal{V}^*$, $v\otimes\varphi=(v\otimes\eta)S(e\otimes\varphi)\in\mathcal{J}$, so $\mathcal{J}$ contains all rank-one maps, hence all finite-rank maps.
\item If $\dim\mathcal{V}<\infty$, then $\mathcal{F}(\mathcal{V})=\LL(\mathcal{V})$, so (ii) makes every nonzero two-sided ideal equal to $\LL(\mathcal{V})$.\\
If $\dim\mathcal{V}=\infty$, then $0\neq\mathcal{F}(\mathcal{V})\neq\LL(\mathcal{V})$, so $\LL(\mathcal{V})$ is not simple.
\end{enumerate}
\hlt{Structural Commentary}: Rank-one operators are the atoms controlling both centrality and ideals in $\LL(\mathcal{V})$.\\
\hlt{Key Insight}: Sandwich one nonzero map to obtain every rank-one map.
\end{proof}
\end{theorem}

\begin{definition} \label{def:algebraic-element-minimal-polynomial-algebra} \hlt{Algebraic Elements and Minimal Polynomials}\\
Let $\mathcal{A}$ be an $\mathbb{F}$-algebra and $a\in\mathcal{A}$.
\begin{enumerate}[label=(\roman*)]
\setlength{\itemsep}{0pt}
\item The \hlt{evaluation map at $a$} is $\operatorname{ev}_a:\mathbb{F}[x]\to\mathcal{A}$, $p\mapsto p(a)$.
\item $a$ is \hlt{algebraic over $\mathbb{F}$} iff $\ker\operatorname{ev}_a\neq0$.
\item If $a$ is algebraic, its \hlt{minimal polynomial} $m_a$ is the unique monic generator of $\ker\operatorname{ev}_a$, equivalently the monic polynomial of least degree with $m_a(a)=0$.
\end{enumerate}
\end{definition}

\begin{theorem} \label{thm:finite-dimensional-algebras-algebraic-elements} \hlt{Finite-Dimensional Algebras Are Algebraic}\\
If $\dim_{\mathbb{F}}\mathcal{A}<\infty$, then every $a\in\mathcal{A}$ is algebraic over $\mathbb{F}$. Moreover, the algebraic closure $\overline{\mathbb{Q}}_{\mathbb{C}}$ of $\mathbb{Q}$ inside $\mathbb{C}$ is an infinite-dimensional $\mathbb{Q}$-algebra whose elements are all algebraic.
\begin{proof}
Let $\dim_{\mathbb{F}}\mathcal{A}=n$ and $a\in\mathcal{A}$. The $n+1$ vectors $1,a,\ldots,a^n$ are linearly dependent, so $\sum_{j=0}^{n}c_ja^j=0$ for some $(c_0,\ldots,c_n)\neq0$; hence $p(x)=\sum_{j=0}^{n}c_jx^j\neq0$ satisfies $p(a)=0$.\\
If $\alpha,\beta\in\overline{\mathbb{Q}}_{\mathbb{C}}$, then $\mathbb{Q}(\alpha,\beta)/\mathbb{Q}$ is finite, so $\alpha+\beta$, $\alpha\beta$, and $r\alpha$ for $r\in\mathbb{Q}$ are algebraic; hence $\overline{\mathbb{Q}}_{\mathbb{C}}$ is a $\mathbb{Q}$-algebra.\\
For each $n\geq1$, $x^n-2$ is Eisenstein at $2$, so it has a root of degree $n$ in $\overline{\mathbb{Q}}_{\mathbb{C}}$; therefore $\dim_{\mathbb{Q}}\overline{\mathbb{Q}}_{\mathbb{C}}=\infty$.\\
\hlt{Structural Commentary}: Finite-dimensionality forces powers of each element into a linear relation.\\
\hlt{Key Insight}: Algebraic need not mean finite-dimensional as an algebra.
\end{proof}
\end{theorem}

\begin{proposition} \label{prop:generated-algebra-by-one-algebraic-element} \hlt{The Algebra Generated by One Algebraic Element}\\
Let $a\in\mathcal{A}$ be algebraic with minimal polynomial $m_a$. Then $\mathbb{F}[a]\defeq\operatorname{im}\operatorname{ev}_a=\{p(a):\deg p<\deg m_a\}$ and $\mathbb{F}[a]\cong\mathbb{F}[x]/\langle m_a\rangle$.
\begin{proof}
$\operatorname{ev}_a:\mathbb{F}[x]\to\mathcal{A}$ is an $\mathbb{F}$-algebra homomorphism with image $\mathbb{F}[a]$.\\
Since $\ker\operatorname{ev}_a=\{p:p(a)=0\}$ and $m_a$ is the monic generator of this kernel, $\ker\operatorname{ev}_a=\langle m_a\rangle$.\\
Proposition \ref{prop:algebra-kernel-first-isomorphism} gives $\mathbb{F}[a]\cong\mathbb{F}[x]/\langle m_a\rangle$. $\forall p\in\mathbb{F}[x]$, by division $p=qm_a+r$, $\deg r<\deg m_a$, hence $p(a)=r(a)$.\\
\hlt{Structural Commentary}: One algebraic element gives one polynomial quotient.\\
\hlt{Key Insight}: Evaluation turns polynomial division modulo $m_a$ into normal forms in $\mathbb{F}[a]$.
\end{proof}
\end{proposition}

\begin{theorem} \label{thm:algebraic-element-annihilator-and-invertibility} \hlt{Minimal Polynomial, Annihilator, and Invertibility}\\
Let $a\in\mathcal{A}$ be algebraic and write $m_a(x)=x^n+\sum_{j=0}^{n-1}\alpha_jx^j$.\\
Then $\operatorname{ann}_{\mathbb{F}[x]}(a)\defeq\{f\in\mathbb{F}[x]:f(a)=0\}=\langle m_a\rangle$, and $a\in\mathcal{A}^{\times}\Longleftrightarrow\alpha_0\neq0$.
\begin{proof} By Definition \ref{def:algebraic-element-minimal-polynomial-algebra}, $\operatorname{ann}_{\mathbb{F}[x]}(a)=\ker\operatorname{ev}_a=\langle m_a\rangle$.
\begin{enumerate}[label=\roman*.]
\setlength{\itemsep}{0pt}
\item[$(\Rightarrow)$] If $a\in\mathcal{A}^{\times}$ and $\alpha_0=0$, then $m_a=xq$ with $\deg q<\deg m_a$, and $0=m_a(a)=aq(a)$ gives $q(a)=0$,\\
contradicting minimality of $m_a$; hence $\alpha_0\neq0$.
\item[$(\Leftarrow)$] If $\alpha_0\neq0$, then $0=m_a(a)=a(a^{n-1}+\sum_{j=1}^{n-1}\alpha_ja^{j-1})+\alpha_0$, so $b\defeq-\alpha_0^{-1}(a^{n-1}+\sum_{j=1}^{n-1}\alpha_ja^{j-1})$ satisfies $ab=1_{\mathcal{A}}$. Since $b\in\mathbb{F}[a]$, $ab=ba$, so $ba=1_{\mathcal{A}}$ and $a\in\mathcal{A}^{\times}$.
\end{enumerate}
\hlt{Structural Commentary}: Invertibility is read from the constant term of the minimal polynomial.\\
\hlt{Key Insight}: A nonzero constant term solves for $a^{-1}$ as a polynomial in $a$.
\end{proof}
\end{theorem}

\begin{definition} \label{def:algebra-zero-divisors-and-spectrum} \hlt{Zero Divisors and Spectrum}\\
Let $\mathcal{A}$ be an $\mathbb{F}$-algebra and $0\neq a\in\mathcal{A}$.
\begin{enumerate}[label=(\roman*)]
\setlength{\itemsep}{0pt}
\item $a$ is a \hlt{left zero divisor} iff $\exists b\in\mathcal{A}\setminus\{0\}$ with $ab=0$.
\item $a$ is a \hlt{right zero divisor} iff $\exists c\in\mathcal{A}\setminus\{0\}$ with $ca=0$.
\item If $a$ is algebraic, \hlt{spectrum} is $\operatorname{spec}_{\mathcal{A}}(a)\defeq\{r\in\mathbb{F}:m_a(r)=0\}$; elements are algebraic \hlt{eigenvalues} of $a$.
\end{enumerate}
\end{definition}

\begin{theorem} \label{thm:algebraic-invertible-iff-not-zero-divisor} \hlt{Invertible if and only if Not a Zero Divisor}\\
Let $a\in\mathcal{A}$ be algebraic. Then $a\in\mathcal{A}^{\times}$ iff $a$ is neither a left nor a right zero divisor.
\begin{proof}
\begin{enumerate}[label=\roman*.]
\setlength{\itemsep}{0pt}
\item[$(\Rightarrow)$] If $a\in\mathcal{A}^{\times}$ and $ab=0$, then $b=a^{-1}ab=0$; if $ca=0$, then $c=caa^{-1}=0$.\\
Hence $a$ is neither a left nor a right zero divisor.
\item[$(\Leftarrow)$] Prove the contrapositive. If $a\notin\mathcal{A}^{\times}$, Theorem \ref{thm:algebraic-element-annihilator-and-invertibility} gives $m_a(x)=xp(x)$ with $p\neq0$ and $\deg p<\deg m_a$.\\
Then $0=m_a(a)=ap(a)=p(a)a$, while minimality gives $p(a)\neq0$. Thus $a$ is both left, right zero divisor.
\end{enumerate}
\hlt{Structural Commentary}: For algebraic elements, noninvertibility forces a polynomial kernel witness.\\
\hlt{Key Insight}: Factoring $m_a$ by $x$ manufactures the zero-divisor witness $p(a)$.
\end{proof}
\end{theorem}

\begin{corollary} \label{cor:algebraic-one-sided-inverse-two-sided} \hlt{One-Sided Inverses Become Two-Sided}\\
Let $a\in\mathcal{A}$ be algebraic. If $\exists b\in\mathcal{A}$ with $ba=1_{\mathcal{A}}$ or $ab=1_{\mathcal{A}}$, then $a\in\mathcal{A}^{\times}$.
\begin{proof}
Assume $ba=1_{\mathcal{A}}$. If $ac=0$, then $c=bac=0$; if $ax=0$, then $x=bax=0$. Hence $a$ is neither a right nor a left zero divisor, so Theorem \ref{thm:algebraic-invertible-iff-not-zero-divisor} gives $a\in\mathcal{A}^{\times}$. The case $ab=1_{\mathcal{A}}$ is symmetric.\\
\hlt{Structural Commentary}: Algebraicity upgrades one-sided cancellation to two-sided invertibility.\\
\hlt{Key Insight}: A one-sided inverse rules out both zero-divisor obstructions.
\end{proof}
\end{corollary}

\begin{theorem} \label{thm:minimal-polynomial-roots-by-algebra-invertibility} \hlt{Roots Detected by Invertibility}\\
Let $a\in\mathcal{A}$ be algebraic and $r\in\mathbb{F}$. Then $m_a(r)=0\Longleftrightarrow a-r1_{\mathcal{A}}\notin\mathcal{A}^{\times}$. In particular, $a\in\mathcal{A}^{\times}\Longleftrightarrow0\notin\operatorname{spec}_{\mathcal{A}}(a)$.
\begin{proof}
\begin{enumerate}[label=\roman*.]
\setlength{\itemsep}{0pt}
\item[$(\Rightarrow)$] If $m_a(r)=0$, write $m_a=(x-r)q$ with $\deg q<\deg m_a$.\\
Then $0=m_a(a)=(a-r1_{\mathcal{A}})q(a)=q(a)(a-r1_{\mathcal{A}})$ and $q(a)\neq0$ by minimality, so $a-r1_{\mathcal{A}}\notin\mathcal{A}^{\times}$.
\item[$(\Leftarrow)$] If $m_a(r)\neq0$, then $\gcd(x-r,m_a)=1$, so $\exists u,v\in\mathbb{F}[x]$ with $u(x)(x-r)+v(x)m_a(x)=1$.\\
Evaluating at $a$ gives $u(a)(a-r1_{\mathcal{A}})=1_{\mathcal{A}}=(a-r1_{\mathcal{A}})u(a)$, hence $a-r1_{\mathcal{A}}\in\mathcal{A}^{\times}$, contradiction.
\end{enumerate}
Taking $r=0$ gives $a\in\mathcal{A}^{\times}\Longleftrightarrow m_a(0)\neq0\Longleftrightarrow0\notin\operatorname{spec}_{\mathcal{A}}(a)$.\\
\hlt{Structural Commentary}: Roots of $m_a$ are precisely the nonunit scalar shifts of $a$.\\
\hlt{Key Insight}: Bézout turns $\gcd(x-r,m_a)=1$ into an explicit inverse.
\end{proof}
\end{theorem}

\begin{theorem} \label{thm:algebra-polynomial-spectral-mapping} \hlt{Polynomial Spectral Mapping for Algebraic Elements}\\
Let $\mathbb{F}$ be algebraically closed, $a\in\mathcal{A}$ algebraic, and $p\in\mathbb{F}[x]$.\\
Then $p(a)$ is algebraic and $\operatorname{spec}_{\mathcal{A}}(p(a))=p(\operatorname{spec}_{\mathcal{A}}(a))\defeq\{p(r):r\in\operatorname{spec}_{\mathcal{A}}(a)\}$.
\begin{proof}
Since $\mathbb{F}[p(a)]\subseteq\mathbb{F}[a]$ and $\dim_{\mathbb{F}}\mathbb{F}[a]<\infty$, $p(a)$ is algebraic. If $p$ is constant, the claim is immediate.\\
If $\lambda\in\operatorname{spec}_{\mathcal{A}}(p(a))$, then $p(a)-\lambda1_{\mathcal{A}}\notin\mathcal{A}^{\times}$. Factor $p(x)-\lambda=c\prod_{j=1}^{m}(x-r_j)^{e_j}$ over $\mathbb{F}$.\\
Since a product of units is a unit, $a-r_j1_{\mathcal{A}}\notin\mathcal{A}^{\times}$ for some $j$. Theorem \ref{thm:minimal-polynomial-roots-by-algebra-invertibility} gives $r_j\in\operatorname{spec}_{\mathcal{A}}(a)$ and $\lambda=p(r_j)$.\\
Conversely, if $r\in\operatorname{spec}_{\mathcal{A}}(a)$, then $p(x)-p(r)=(x-r)q(x)$, so $p(a)-p(r)1_{\mathcal{A}}=(a-r1_{\mathcal{A}})q(a)$ is noninvertible; otherwise $a-r1_{\mathcal{A}}$ has a one-sided inverse and Corollary \ref{cor:algebraic-one-sided-inverse-two-sided} makes it invertible. Thus $p(r)\in\operatorname{spec}_{\mathcal{A}}(p(a))$.\\
\hlt{Structural Commentary}: Polynomial evaluation carries spectra by applying the same polynomial to roots.\\
\hlt{Key Insight}: Factor $p(x)-\lambda$ and detect a nonunit scalar shift of $a$.
\end{proof}
\end{theorem}

\begin{proposition} \label{prop:algebra-spectra-of-ab-and-ba} \hlt{Spectra of $ab$ and $ba$}\\
If $\dim_{\mathbb{F}}\mathcal{A}<\infty$, $\mathbb{F}$ is algebraically closed, and $a,b\in\mathcal{A}$, then $\operatorname{spec}_{\mathcal{A}}(ab)=\operatorname{spec}_{\mathcal{A}}(ba)$.
\begin{proof}
All elements are algebraic by Theorem \ref{thm:finite-dimensional-algebras-algebraic-elements}. Let $r\in\mathbb{F}^{\times}$.\\
If $ba-r1_{\mathcal{A}}\in\mathcal{A}^{\times}$, set $u\defeq a(ba-r1_{\mathcal{A}})^{-1}b-1_{\mathcal{A}}$; then $(ab-r1_{\mathcal{A}})u=r1_{\mathcal{A}}$, so $ab-r1_{\mathcal{A}}$ has a right inverse, hence is invertible by Corollary \ref{cor:algebraic-one-sided-inverse-two-sided}. Thus $r\notin\operatorname{spec}_{\mathcal{A}}(ba)\Rightarrow r\notin\operatorname{spec}_{\mathcal{A}}(ab)$, and symmetry gives equality off $0$.\\
If $ba\in\mathcal{A}^{\times}$ with inverse $y$, then $b(ay)=1_{\mathcal{A}}$ and $(yb)a=1_{\mathcal{A}}$, so $b,a\in\mathcal{A}^{\times}$ by Corollary \ref{cor:algebraic-one-sided-inverse-two-sided}; hence $ab\in\mathcal{A}^{\times}$. The converse is symmetric, so $0$ belongs to both spectra or neither.\\
\hlt{Structural Commentary}: $\operatorname{spec}(ab)\setminus\{0\}=\operatorname{spec}(ba)\setminus\{0\}$, and algebraicity resolves the $0$-case.\\
\hlt{Key Insight}: The resolvent identity transfers invertibility between $ba-r1_{\mathcal{A}}$ and $ab-r1_{\mathcal{A}}$ for $r\neq0$.
\end{proof}
\end{proposition}

\begin{definition} \label{def:division-algebra} \hlt{Division Algebras}\\
An associative $\mathbb{F}$-algebra $\mathcal{D}$ is a \hlt{division algebra} iff $\mathcal{D}\setminus\{0\}\subseteq\mathcal{D}^{\times}$.
\end{definition}

\begin{example} \label{ex:hamilton-quaternions} \hlt{Hamilton's Quaternions}\\
Hamilton's quaternion algebra is $\mathbb{H}\defeq\mathbb{R}1\oplus\mathbb{R}i\oplus\mathbb{R}j\oplus\mathbb{R}k$ with unit $1$,\\
relations $i^2=j^2=k^2=-1$, $ij=k=-ji$, $jk=i=-kj$, and $ki=j=-ik$, extended $\mathbb{R}$-bilinearly.\\
It is noncommutative since $ij=-ji$. For $q=r_0+r_1i+r_2j+r_3k$, define $\overline{q}\defeq r_0-r_1i-r_2j-r_3k$ and $N(q)\defeq q\overline{q}=r_0^2+r_1^2+r_2^2+r_3^2$.
If $q\neq0$, then $N(q)>0$ and $q^{-1}=N(q)^{-1}\overline{q}$, so $\mathbb{H}$ is a real division algebra.
\end{example}

\begin{figure}[H]
  \centering
  \begin{tikzpicture}[
  >=Stealth,
  every node/.style={font=\small},
  basis/.style={circle, draw=blue!60!black, fill=blue!7,
    minimum size=0.78cm, inner sep=1pt},
  fwdedge/.style={->, line width=0.9pt, draw=dartmouthgreen!70!black},
  lbl/.style={font=\scriptsize, fill=white, inner sep=1.5pt}
]
  \node[basis] (i) at (90:1.55) {$i$};
  \node[basis] (j) at (210:1.55) {$j$};
  \node[basis] (k) at (330:1.55) {$k$};

  \draw[fwdedge] (i) to[bend left=18] node[lbl, left=2pt] {$ij=k$} (j);
  \draw[fwdedge] (j) to[bend left=18] node[lbl, below=3pt] {$jk=i$} (k);
  \draw[fwdedge] (k) to[bend left=18] node[lbl, right=2pt] {$ki=j$} (i);
\end{tikzpicture}

  \caption{Quaternion basis products: the cycle $i\to j\to k\to i$ gives
  $ij=k$, $jk=i$, and $ki=j$; reversed order negates the product. Also
  $i^2=j^2=k^2=-1$, and $1$ is the unit.}
  \label{fig:quaternion-multiplication-cycle}
\end{figure}

\begin{theorem} \label{thm:division-algebra-algebraically-closed-field} \hlt{Division Algebras over Algebraically Closed Fields}\\
If $\dim_{\mathbb{F}}\mathcal{D}<\infty$, $\mathcal{D}$ is a division $\mathbb{F}$-algebra, and $\mathbb{F}$ is algebraically closed, then $\mathcal{D}=\mathbb{F}1_{\mathcal{D}}\cong\mathbb{F}$.
\begin{proof}
Let $a\in\mathcal{D}$. By Theorem \ref{thm:finite-dimensional-algebras-algebraic-elements}, $a$ is algebraic with minimal polynomial $m_a$.\\
If $m_a=pq$ with $0<\deg p,\deg q<\deg m_a$, then $0=m_a(a)=p(a)q(a)$, while minimality gives $p(a),q(a)\neq0$; this contradicts $\mathcal{D}\setminus\{0\}\subseteq\mathcal{D}^{\times}$. Thus $m_a$ is irreducible, hence linear over algebraically closed $\mathbb{F}$: $m_a=x-r$.\\
Therefore $a=r1_{\mathcal{D}}$, so $\mathcal{D}\subseteq\mathbb{F}1_{\mathcal{D}}$, and the reverse inclusion is the scalar copy.\\
\hlt{Structural Commentary}: Algebraic closure leaves only scalars.\\
\hlt{Key Insight}: Irreducible minimal polynomials collapse to linear polynomials.
\end{proof}
\end{theorem}

\begin{theorem} \label{thm:wedderburn-finite-division-algebra} \hlt{Wedderburn's Theorem}\\
Every finite division algebra $\mathcal{D}$ is commutative; equivalently $\mathcal{D}=Z(\mathcal{D})$, hence $\mathcal{D}$ is a finite field.
\begin{proof}
Let $F\defeq Z(\mathcal{D})$, $|F|=z$, $[\mathcal{D}:F]=n$, and $G\defeq\mathcal{D}^{\times}$.\\
Suppose $\mathcal{D}$ is noncommutative, and let $B$ be representatives of the noncentral conjugacy classes of $G$.\\
For $\beta\in B$, $C_{\mathcal{D}}(\beta)$ is a proper division subalgebra containing $F$, so $|C_{\mathcal{D}}(\beta)|=z^{k_\beta}$ with $1<k_\beta<n$ and $k_\beta\mid n$.\\
By the class equation, Theorem \ref{thm:finite-group-class-equation}, $z^n-1=z-1+\sum_{\beta\in B}(z^n-1)/(z^{k_\beta}-1)$.\\
Since $k_\beta\mid n$ and $k_\beta<n$, the primitive $n$th roots occur in $(x^n-1)/(x^{k_\beta}-1)$, so Propositions \ref{prop:roots-of-unity-cyclic-and-cyclotomic-factorization} and \ref{prop:cyclotomic-polynomials-integer-coefficients} give $Q_n(z)\mid(z^n-1)/(z^{k_\beta}-1)$ for each $\beta$, and also $Q_n(z)\mid z^n-1$.\\
The class equation therefore gives $Q_n(z)\mid z-1$, contradicting Lemma \ref{lem:cyclotomic-size-estimate-for-wedderburn}, which gives $|Q_n(z)|>z-1$.\\
Thus $G$ has no noncentral conjugacy classes, so $G$ is abelian and $\mathcal{D}$ is commutative.\\
\hlt{Structural Commentary}: The class equation converts noncommutativity into cyclotomic divisibility.\\
\hlt{Key Insight}: The cyclotomic lower bound is too large to divide $z-1$.
\end{proof}
\end{theorem}

\begin{theorem} \label{thm:frobenius-real-division-algebras} \hlt{Frobenius Classification over $\mathbb{R}$}\\
If $\dim_{\mathbb{R}}\mathcal{D}<\infty$ and $\mathcal{D}$ is a real division algebra, then $\mathcal{D}\cong\mathbb{R}$, $\mathcal{D}\cong\mathbb{C}$, or $\mathcal{D}\cong\mathbb{H}$.
\begin{proof}
For each $d\in\mathcal{D}$, $m_d$ is linear or irreducible quadratic over $\mathbb{R}$; completing the square gives $d=r+u$ with $r\in\mathbb{R}$ and $u^2\in\mathbb{R}_{\le0}$. Set $\mathcal{D}_0\defeq\{u\in\mathcal{D}:u^2\in\mathbb{R}_{\le0}\}$; the preceding decomposition and uniqueness give $\mathcal{D}=\mathbb{R}\oplus\mathcal{D}_0$. If $\mathcal{D}_0=0$, then $\mathcal{D}\cong\mathbb{R}$. If $\dim\mathcal{D}_0=1$, choose $i\in\mathcal{D}_0$ with $i^2=-1$, so $\mathcal{D}=\mathbb{R}\oplus\mathbb{R}i\cong\mathbb{C}$.\\
If $\dim\mathcal{D}_0>1$, use $\langle u,v\rangle\defeq-\frac12(uv+vu)$ on $\mathcal{D}_0$ to choose orthonormal $i,j\in\mathcal{D}_0$ with $i^2=j^2=-1$; $ij+ji=0$.\\
Then $k\defeq ij$ satisfies $k^2=-1$, $ij=k$, $jk=i$, $ki=j$, and reversed products change sign.\\
If $\ell\in\mathcal{D}_0$ is orthogonal to $i,j,k$, the same identities force $\ell=0$, so $\mathcal{D}=\mathbb{R}\oplus\mathbb{R}i\oplus\mathbb{R}j\oplus\mathbb{R}k\cong\mathbb{H}$.\\
\hlt{Structural Commentary}: A real division algebra splits into scalars plus a Euclidean pure part.\\
\hlt{Key Insight}: Orthogonal imaginary units generate exactly $\mathbb{C}$ or $\mathbb{H}$.
\end{proof}
\end{theorem}

\begin{perspective} \label{per:algebra-chapter-operator-connection} \hlt{Operator Connection}\\
$T\in\LL(\mathcal{V})$ gives an $\mathbb{F}[x]$-module by $x\cdot v=T(v)$. For any algebra $\mathcal{A}$, the left regular map $L_{\mathcal{A}}:\mathcal{A}\hookrightarrow\LL(\mathcal{A})$ realizes $\mathcal{A}$ as a faithful endomorphism algebra.
\end{perspective}
